%% kapa: exam=1 round=30 year=2019 period=2 form=A subject=appraisal_law subject_ko=감정평가관계법규 seq=1 number=1
%% source: Q-Net 감정평가사 기출문제 (cst00309) artlSeq=5205205 file=extracted/1008/5205205/2019년 제30회 감정평가사 1차시험 문제지/2019년 제30회 감정평가사 1차시험 문제지/2019년도 제30회 감정평가사 1차 2교시 A형.hwp
%% conversion: kordoc 4.24.51 -> tools/gen1.py (auto); review=proofread-sonnet; flags=none
\begin{question}{1}
국토의 계획 및 이용에 관한 법령상 도시․군관리계획에 관한 설명으로 옳지 않은 것은?
\begin{choices}
  \item 도시ㆍ군관리계획 결정의 효력은 그 결정이 있는 때부터 발생한다.
  \item 도시ㆍ군관리계획은 광역도시계획과 도시ㆍ군기본계획에 부합되어야 한다.
  \item 수산자원보호구역의 지정에 관한 도시ㆍ군관리계획은 해양수산부장관이 결정할 수 있다.
  \item 주민은 도시ㆍ군관리계획을 입안할 수 있는 자에게 기반시설의 개량에 관한 사항에 대하여 도시ㆍ군관리계획의 입안을 제안할 수 있다.
  \item 국토교통부장관이 직접 지형도면을 작성한 경우에는 이를 고시하여야 한다.
\end{choices}
\end{question}
\answer{1}
